\subsection{Sprint 4}

Para esta sprint, a equipe havia previsto entregar a tela dos dashboards no
front-end, bem como a parte da presença digital, conciliando todo o conhecimento
obtido das spikes passadas, e o deploy na AWS do aplicativo. Da minha parte,
fiquei com o front-end, tendo em vista que, de imediato, era a função na qual eu
mais conseguiria colaborar com o time, já que as tasks do back-end com FastAPI,
tecnologia que estudei, já haviam sido finalizadas.

Da parte da presença digital, a equipe conseguiu entregar um resultado bem
satisfatório, principalmente na visão do stakeholder, que ficou feliz com a prova de
conceito apresentada. O deploy foi feito conforme esperado, no entanto a integração
das EC2 estava apresentando falha. Já o front-end foi onde estive mais presente
nesta sprint: iniciei a tela de dashboards, adicionando os dois primeiros gráficos e
estruturando o código, deixando os outros dois para meus colegas.

Apesar de uma boa entrega, a sprint foi conturbada. Houve problemas comuns,
como provas e trabalhos em paralelo, que apesar de já não preocuparem tanto,
ainda eram um fator de complicação a ser levado em conta. Além disso, houve
problemas mais técnicos, como memory leaks descobertos um dia antes da
apresentação e problemas de integração no deploy. Por fim, a API utilizada para a
pesquisa da presença digital tinha um limite de requisições muito baixo,
ocasionando problemas ao ser chamada muitas vezes no mesmo dia.

Mediante os problemas, é possível concluir que a AGES requer atenção maior do
que outras cadeiras, até pela carga horária excedente que, apesar de quatro
créditos, consome algo equivalente a mais horários semanais. Além disso, ficou
como aprendizado sempre deixar tempo livre antes e depois do deploy, pois ao
redor desse horário é comum surgirem problemas e, quanto maior a importância
para o time, mais crucial é a presença imediata para resolvê-los. Pensando em uma
futura AGES III, isso reforça a importância de ter atenção ao deploy e trabalhar com
code freezes prévios aos deadlines.

Por fim, o problema da API não tinha uma solução ótima, tendo em vista que ela
era a única opção disponível para aquela função e satisfazia a demanda do cliente.
Para a prova de conceito, isso ficou como uma limitação, mas ao assumir um
investimento futuro do cliente na API do Google, esse problema deixaria de ser um
bloqueio.

Após a apresentação para o cliente, recebemos bons feedbacks, ressaltamos
pontos importantes a serem mudados e levantamos ideias do que corrigir, além de
adicionar mais testes no tempo restante. Com mais uma semana até a apresentação
geral dos projetos, planejamos encorpar um pouco mais o projeto para entregar algo
mais polido.
